\documentclass[11pt,a4paper]{article}
\usepackage[utf8]{inputenc}
\usepackage[T1]{fontenc}
\usepackage[brazil]{babel}
\usepackage[a4paper,margin=2.3cm]{geometry}
\usepackage{enumitem}
\usepackage{hyperref}
\usepackage{xcolor}
\setlength{\parindent}{0pt}
\setlength{\parskip}{5pt}
\sloppy
\setlength{\emergencystretch}{3em}
\hypersetup{colorlinks=true,urlcolor=blue,linkcolor=black}
\newcommand{\cabecalho}[2]{\begin{center}{\small \textbf{Sprint 2 -- #1}}\\[1mm]\textit{Inventário de Ativos e Vulnerabilidades}\\Matheus Moreli Durão Silva -- 12621CBS217\\Cibersegurança -- Bacharelado -- Noturno\end{center}\hrule\vspace{4mm}}
\newcommand{\bloco}[1]{\par\vspace{2mm}\textbf{#1}\par}
\newcommand{\codigo}[1]{\texttt{#1}}

\begin{document}\cabecalho{S2\_12 -- Fundamentos de Vulnerabilidades}{}
\bloco{Conceitos} Vulnerabilidade é uma falha ou fraqueza que pode ser explorada. Ameaça é algo ou alguém capaz de explorar uma vulnerabilidade. Risco é a combinação da possibilidade de exploração com o impacto no contexto do ativo. Não são sinônimos.
\bloco{Dados verificáveis} CVE identifica uma vulnerabilidade divulgada publicamente; CWE classifica a fraqueza; CVSS registra uma pontuação técnica. O sistema também guarda fonte e data, ativo afetado, impacto, prioridade, tratamento e como o tratamento será verificado. A prioridade considera o contexto: uma mesma nota CVSS pode ser mais urgente em um ativo crítico.
\bloco{Fontes de consulta} O programa CVE cataloga vulnerabilidades divulgadas; a NVD reúne dados de gestão e métricas; a CWE descreve fraquezas de software e hardware. Consultar a referência antes do cadastro impede apresentar suposição como fato.\newline CVE: \url{https://www.cve.org/}; NVD: \url{https://nvd.nist.gov/}; CWE: \url{https://cwe.mitre.org/}.
\bloco{Caso real 1: CVE-2019-0708} A NVD descreve execução remota de código no Remote Desktop Services de versões antigas do Windows e classifica a fraqueza como CWE-416. Antes de associar ao notebook de demonstração, é necessário verificar versão do sistema e exposição do RDP. Tratamento: correção do fabricante e restrição do acesso remoto. Fonte: \url{https://nvd.nist.gov/vuln/detail/CVE-2019-0708}.
\bloco{Caso real 2: CVE-2021-44228} A NVD registra falha no Log4j2 associada a JNDI e execução remota de código em configurações afetadas. O servidor de demonstração só estaria exposto se a aplicação usasse versão afetada de \codigo{log4j-core}. Tratamento: atualizar a biblioteca e dependências; verificar componentes instalados. Fonte: \url{https://nvd.nist.gov/vuln/detail/CVE-2021-44228}.
\bloco{Limite da análise} Os ativos do JSON são fictícios. CVE real não prova que um equipamento real usa versão vulnerável. O estudo detalhado, incluindo impacto, tratamento e fonte primária do fornecedor, está em \codigo{docs/cves-analisados.md} no repositório principal.
\end{document}
